\subsection{Thesauri and WordNet}
\label{sec:wordnet}
\slidesrc{Semantics (long version): 22--30}

\subsubsection{Structured lexicons and thesauri}
\slidesrc{Semantics (long version): 22--23}

A \textbf{thesaurus} (structured lexicon) is an alternative to the alphabetical
dictionary: a list of words grouped by \emph{meaning}, with a \textbf{hierarchical}
organisation, like the taxonomies of the natural sciences (children are ``types of'' the
father and inherit its properties). For ordinary words the hierarchy is given by hyponymy
and synonymy (Section~\ref{sec:lexical-relations}).
\begin{itemize}[nosep]
  \item It shows the structure of lexical knowledge, and is much used for technical
        terminology.
  \item It can be enriched with other relations: antonyms, part/member relations
        (\emph{has-as-part}, \emph{has-as-member}), not just \emph{is-a-type-of}.
  \item It can be explored in any direction (up, down, across), and obvious
        \textbf{distance metrics} (e.g.\ the number of links between two words) measure
        their similarity.
\end{itemize}

\subsubsection{WordNet}
\slidesrc{Semantics (long version): 24--30}

\textbf{WordNet} is a large lexical database of English, started in \textbf{1985} at
Princeton University by a group of psychologists and linguists (led by George Miller,
the author of the quotation in Section~\ref{sec:distributional-hypothesis}). Its
theoretical basis is psycholinguistics: what are the properties of the \emph{mental
lexicon}? Global organisation: \textbf{nouns} are organised as topical hierarchies,
\textbf{verbs} by entailment relations.

\begin{nbdef}{Synset}
In WordNet the basic unit is the \textbf{synset} (synonym set). A word meaning is
represented simply by listing the word forms that can express it. Synsets are
unambiguous designators: the two main senses of \emph{board} are
\begin{center}
\{board, plank, \ldots\} \quad (a piece of lumber) \qquad vs.\qquad \{board, committee, \ldots\}
\quad (a group of people assembled for some purpose).
\end{center}
\end{nbdef}

\begin{itemize}[nosep]
  \item Members of a synset are rarely true synonyms: WordNet does not try to capture
        subtle distinctions (connotation, collocation) among them.
  \item Synsets are often enough to \emph{differentiate} senses. If no suitable synonym
        exists, a short \textbf{gloss} is used, e.g.\ \{board, (a person's meals,
        provided regularly for money)\}. WordNet gives a gloss for every meaning, and
        often an example sentence.
\end{itemize}

\begin{nbexample}{the noun senses of \emph{board} in WordNet (excerpt)}
\small
\begin{itemize}[nosep]
  \item board (a committee having supervisory powers) ``\emph{the board has seven
        members}'';
  \item board, plank (a stout length of sawn timber; made in a wide variety of sizes
        and used for many purposes);
  \item board (a flat piece of material designed for a special purpose) ``\emph{he
        nailed boards across the windows}'';
  \item board, table (food or meals in general) ``\emph{she sets a fine table}'';
  \item display panel, display board, board (a vertical surface on which information
        can be displayed to public view);
  \item dining table, board (a table at which meals are served);
  \item control panel, instrument panel, control board, board, panel (electrical device
        consisting of a flat insulated surface that contains switches and dials).
\end{itemize}
For the first sense WordNet also lists its relations: its members (\emph{director},
\emph{board member}: meronyms), its hypernym (\emph{committee}), its sister terms.
\end{nbexample}

\paragraph{Size.} WordNet 3.0:
\begin{center}
\footnotesize
\setlength{\tabcolsep}{4pt}
\begin{tabular}{@{}l rrr @{\hspace{0.6cm}} rrr@{}}
\toprule
 & \textbf{Words} & \textbf{Synsets} & \textbf{Word--sense} & \textbf{Monos.} & \textbf{Polys.} & \textbf{Polys.} \\
 &  &  & \textbf{pairs} & \textbf{words} & \textbf{words} & \textbf{senses} \\
\midrule
Noun      & 117\,798 & 82\,115  & 146\,312 & 101\,863 & 15\,935 & 44\,449 \\
Verb      & 11\,529  & 13\,767  & 25\,047  & 6\,277   & 5\,252  & 18\,770 \\
Adjective & 21\,479  & 18\,156  & 30\,002  & 16\,503  & 4\,976  & 14\,399 \\
Adverb    & 4\,481   & 3\,621   & 5\,580   & 3\,748   & 733     & 1\,832 \\
\midrule
Total     & 155\,287 & 117\,659 & 206\,941 & 128\,391 & 26\,896 & 79\,450 \\
\bottomrule
\end{tabular}
\end{center}
Most words have a single sense (128\,391 of 155\,287), but the polysemous ones are the
frequent ones: verbs in particular have on average $18\,770/5\,252 \approx 3.6$ senses
each.

\begin{center}
\small
\begin{tabularx}{\linewidth}{@{}l >{\raggedright\arraybackslash}X@{}}
\toprule
\textbf{POS} & \textbf{Lexical relations in WordNet} \\
\midrule
Nouns & synonym $\sim$ antonym (opposite of); hypernym (is a kind of) $\sim$ hyponym (for
  example); coordinate (sister) terms: share the same hypernym; holonym (is part of)
  $\sim$ meronym (has as part) \\
Verbs & synonym $\sim$ antonym; hypernym $\sim$ troponym (\emph{lisp} -- \emph{talk});
  entailment (\emph{snore} -- \emph{sleep}: if you snore you sleep); coordinate terms \\
Adjectives, adverbs & the above, plus related nouns, verb participles, derivational
  information \\
\bottomrule
\end{tabularx}
\end{center}

WordNet is available in NLTK (Section~\ref{sec:word-segmentation}) and online
(``WordNet Search''); for Italian and many other languages, multilingual resources such as
BabelNet extend it.
